\section{参考文献}
\renewcommand{\refname}{}% 标题已由上一行给出，避免与宏包自动标题重复
\begin{thebibliography}{99}
\bibitem{fpp3} Hyndman R J, Athanasopoulos G. Forecasting: Principles and Practice[M/OL]. 3rd ed. Melbourne: OTexts, 2021. https://otexts.com/fpp3/（访问时间: 2026-09-13）. 本文引用第 2 章（时间序列图形与自相关）、第 5 章（预测工具箱: 5.2 简单预测方法、5.4 点预测精度评估、5.6 时间序列交叉验证）、第 7 章（时间序列回归模型）.
\bibitem{arrow1951} Arrow K J, Harris T, Marschak J. Optimal inventory policy[J]. Econometrica, 1951, 19(3): 250-272.
\bibitem{glahn1972} Glahn H R, Lowry D A. The use of Model Output Statistics (MOS) in objective weather forecasting[J]. Journal of Applied Meteorology, 1972, 11(8): 1203-1211.
\bibitem{bertsimas2013} Bertsimas D, Litvinov E, Sun X A, et al. Adaptive robust optimization for the security constrained unit commitment problem[J]. IEEE Transactions on Power Systems, 2013, 28(1): 52-63.
\bibitem{mayne2000} Mayne D Q, Rawlings J B, Rao C V, et al. Constrained model predictive control: Stability and optimality[J]. Automatica, 2000, 36(6): 789-814.
\bibitem{hong2016} Hong T, Pinson P, Fan S, et al. Probabilistic energy forecasting: Global Energy Forecasting Competition 2014 and beyond[J]. International Journal of Forecasting, 2016, 32(3): 896-913.
\end{thebibliography}
